\section{引言}
\label{sec:introduction}

CLIP 为零样本异常定位提供了直接的语义接口：视觉 patch 可以与正常和异常状态的语言原型进行匹配。WinCLIP 使用组合式正常/异常 prompt 和局部窗口评分实现零样本或少样本异常分割~\cite{jeong2023winclip}。AnomalyCLIP 从辅助异常数据中学习 object-agnostic 的正常和异常 prompt~\cite{zhou2024anomalyclip}；后续方法进一步通过 hybrid learnable prompts~\cite{cao2024adaclip}、visual context prompting~\cite{qu2024vcpclip} 和 anomaly-aware adapters~\cite{ma2025aaclip} 改进这一类模型。这些工作表明，CLIP 能够为异常定位提供有效的正常/异常语义原型。

局部异常证据并不只由语义相似度决定。工业缺陷通常表现为划痕、污染、压痕、裂纹或缺失部件，其中许多缺陷会引起纹理、边界连续性或细粒度结构的局部变化。相似的局部变化也可能来自正常材质、重复纹理、物体轮廓和成像噪声。因此，一个 patch 响应是否指向异常，取决于它相对于当前图像正常外观的偏离程度。同样的高响应区域，在光滑金属表面上可能是划痕，在编织物或粗糙表面上则可能是正常结构。

这一点揭示了 \clipzsad{} 中一个具体的 reference gap。可学习的正常/异常文本原型定义了跨类别共享的语义轴，但它们没有给出当前测试图像的正常外观参照。对于一张具体图像，表面颗粒度、纹理密度、结构边界和重复图案共同决定了哪些局部变化属于正常范围。缺少这一图像条件化参照时，正常纹理、结构边界和真实缺陷容易在局部响应空间中相互混杂。

文本侧的一个诊断现象进一步指向这一缺口。当训练得到的通用 abnormal prompt feature 之外再加入当前缺陷外观的 oracle detailed prompt 时，部分样例中标注异常区域的响应被进一步激活，异常区域与背景区域的响应间隔也随之增大。该现象说明，通用文本原型之外，当前测试图仍携带能够调节 patch response 的外观信息。本文将这一信息转到图像侧处理，从测试图自身选择可靠 patch 证据，估计图像条件化正常外观基准。

本文提出面向 \clipzsad{} 的图像条件化正常参照估计。给定单张测试图像，\method{} 保持 CLIP 和 AnomalyCLIP 编码器冻结，在 CLIP patch 特征场中选择可靠证据。所选正常证据被聚合为图像条件化正常参照，随后每个 patch 的异常响应相对于该参照和共享的正常/异常语义轴重新计算。整个过程不使用目标类别正常样本库、标签或测试时参数更新。

方法的关键在于 patch 证据选择。如果证据选择只由语义响应驱动，正常纹理、结构边界和缺陷区域仍可能混在一起。\method{} 因此在 CLIP patch 特征网格上测量局部结构变化。频域响应用于刻画这种变化，并约束哪些 patch 可以支撑正常参照。语义异常响应较低且局部结构稳定的 patch 构成正常参照证据；语义可疑且结构变化显著的 patch 从正常参照中排除，并在参照估计后重新评估。

本文贡献如下：
\begin{itemize}
  \item 指出 \clipzsad{} 中的 reference gap：可学习的正常/异常文本原型提供跨类别语义轴，但缺少当前测试图像的正常外观参照。
  \item 提出 \method{}，一种训练无关的图像条件化正常参照估计方法，从单张测试图像中选择可靠 patch 证据，并在冻结 CLIP 语义空间中重新计算局部异常响应。
  \item 将频域响应用于局部结构变化刻画和证据选择，并通过受控机制实验验证其角色：Direct 作为负控，Global 检验固定参照，Sem 作为图像条件化基线，\method{} 在 MVTec AD 和 VisA 上取得最佳受控结果。
\end{itemize}
